\subsection{Local propagation at every supplied coefficient curve}
\label{sec:selmer-local-propagation}

The polarized equation model supplies the actual representation at a
normalized characteristic-prime curve.  We now apply the existing
padded-presentation and incidence theorems to that model, including
their tangent and local-compatibility inputs.

Retain the original ordinary pair over \(B\), the normalized point
over \(A=k[[t]]\), its auxiliary roots and inertia parameters, and
the specified independent odd classes \(\kappa_1,\ldots,\kappa_g\).
Put \(F=k((t))\), \(R=\widehat{B_{\ker f}}\), and
\(\mathfrak m=\mathfrak m_R\).  Assume the source's three conditions
\(\mathrm K_1(h,\kappa,D_0)\), \(\mathrm K_2(h,\kappa)\), and
\(\mathrm K_3(h,\kappa)\), with their original groups, local flags
and skew primitives.  The integral-limit comparison supplies their
bounded Laurent versions.

\begin{proposition}[The full local propagation statement]
\label{prop:selmer-local-propagation}
For every pair of primes \(Z,Q\subset R\), if
\[
 \mathrm{DiscSt}(D_R,Z),\qquad \dim(R/Z)\geq1+D_0,
 \qquad
 \mathrm{DiscSt}(D_R,Q)\ \text{or}\ \mathrm{NonSS}(I,R,Q),
\]
then \(Z\subseteq Q\).
Here \(I\) is the original already-closed equation ideal, and its
base change in \(\mathrm{NonSS}\) is the literal algebraic span.
This supplies exactly the source's \(\mathrm{GalOut}\) contract.
\end{proposition}

\begin{proof}
\emph{The actual coefficient and tangent objects.}
The polarized model gives \(\mathcal A=R\oplus N\) and \(v\),
with \(N^2\subseteq\mathfrak m\), transpose involution and the full
prescribed block span.  Evaluation of the original curve variable
gives the coefficient DVR
\[
 O'=\widehat{O[[s]]_{(\pi)}}\longrightarrow R,
 \qquad O'/(\pi)=F.
\]
Its image contains the original coefficients and \(s\); this uses
\(f(s)=t\) and the surjection \(O\to k\).
With \(\mathfrak r=(N^2)\), set
\[
 H=\pi R+\mathfrak m^2+\mathfrak r,\qquad
 T=\operatorname{Der}_{O'}(R,F)_{\mathfrak r=0}
       =(\mathfrak m/H)^*,\qquad
 L=\operatorname{Hom}_R(N,F).
\]

All infinitesimal tests below have bounded Laurent coordinates.
Indeed, for an \(O'\)-constant map to a finite augmented \(F\)-algebra
with augmentation ideal of exponent \(r\), its restriction to \(B\)
kills \(p^r\), where \(p=\ker f\).  The filtration by \(p^j/p^{j+1}\)
makes \(B/p^r\) finite over \(C=O[[s]]\), since \(B/p=C/(\pi)\).
The image is therefore finite over \(A\).  As
\(\mathfrak m_B=(s)+p\), its image of \(\mathfrak m_B^j\) is
contained in \(t^{j-r+1}\) times that finite \(A\)-module for
\(j\geq r-1\).  Apply this to the finite original \(B\)-lattice
of \(v\).  It gives both bounded denominators and the required
continuity, with no assumption about arbitrary maps into \(F\).

\emph{The original ordinary identities give local corrections.}
Over any finite local augmented \(F\)-algebra, retain the ordered
identity
\[
 \prod_i\bigl(\sigma(g_i)-\chi_i(g_i)\bigr)=0
                              \quad\text{for all }(g_i).
\]
The residue has the prescribed flag with each distinct character
once.  Choose \(z\in F[H_v]\) separating the finitely many character
values; finitely many proper hyperplanes cannot cover an
\(F\)-vector space.  Multilinearity permits this group-algebra input.
The Lagrange projectors
\[
 P_i=\prod_{j\ne i}
      \frac{\sigma(z)-\chi_j(z)}{\chi_i(z)-\chi_j(z)}
\]
give its ordered ordinary flag.  To check the order, let
\(J_i=\ker\chi_i\).  The identity \(J_1\cdots J_nE^n=0\)
says that \(J_2\cdots J_nE^n\) is the first eigenline: it is
contained there and contains it by the product of the other
spectral factors.  Repeat on the quotient.

For two representations and character systems agreeing modulo a
square-zero ideal, their projectors satisfy
\[
 S=\sum_iP_i^1P_i^0=1-Y,\qquad S^{-1}=1+Y.
\]
They match the ordered flags.  Thus
\[
 \sigma_1(g)\sigma_0(g)^{-1}-1
       -\bigl(\sigma_0(g)Y\sigma_0(g)^{-1}-Y\bigr)
\]
is upper triangular in the original flag, with diagonal
\((\chi_i^1-\chi_i^0)/\chi_i^0\).  The matrix \(Y\) stays in the
specified parity subalgebra: all projectors belong to it, and unit
inverses are finite geometric polynomials in nilpotents.  The corrected
cocycle's diagonal vanishes on \(\mu_v\).  The original inertia and polarization
identities supply the remaining Selmer conditions.

\emph{The tangent maps are injective modulo the allowed boundaries.}
Writing \(v=v_0+v_1\) and \(\bar\rho=\rho\circ h\), the actual maps are
\[
 a_\alpha=\alpha(v_0)\bar\rho^{-1},\qquad
 b_\mu=\mu(v_1)\bar\rho^{-1}\qquad(\alpha\in T,\ \mu\in L).
\]
Their dual-number tests are ring homomorphisms, and the preceding
argument makes them bounded even and odd Selmer cocycles.
If \(b_\mu=dY\), extend
\(\mu(v_1)=\bar\rho Y-Y\bar\rho\) to the full block span.
A pure cross-corner entry \(m\in N\) has zero even reduction,
forcing \(\mu(m)=0\).

If \(a_\alpha=dY\), conjugation by \(1+\epsilon Y\) makes the
test representation constant, hence fixes all determinant
coefficients.  Its ordered ordinary characters are fixed as well:
in the \(i\)-th diagonal ordinary identity, separate \(\psi_i\)
from each other character using separate group elements.  The
other factors are units, forcing the remaining one to vanish.
The same argument on the actual distinct Frobenius eigenlines
fixes every auxiliary root \(U_{ji}\).  The original partial-root
equation then kills \(y_{ji}\), since \(h(\tau_j)=1\).
The coefficients and \(s\) are fixed by \(O'\)-linearity.
The exact generator condition and bounded continuity give
constancy on \(B\), then on \(R\); therefore \(\alpha=0\).

An invariant even skew matrix is zero: the full block span makes
it scalar on each block, and \(2\ne0\) kills the skew scalars.
Condition \(\mathrm K_3\) and the odd injection now give generators
\(m_1,\ldots,m_g\) of \(N\) such that
\[
 b_\mu-\sum_j\mu(m_j)\kappa_j
                    \quad\text{is an allowed odd boundary}.
\]
Duality gives these generators modulo \(\mathfrak mN\), and
Nakayama lifts generation.  No freeness of \(N\) is used.

\emph{The second-order functional really lifts.}
Let \(\Phi:N\to F\) be residue-linear over \(O'\) and kill \(HN\).
Then
\[
 \alpha_m(r)=\Phi(rm)-\bar r\Phi(m)\in T,\qquad
 \alpha_{rm}=\bar r\alpha_m.
\]
A basis \(\alpha_k\) of \(T\) gives \(\mu_k\in L\) with
\(\Phi(rm)=\bar r\Phi(m)+\sum_k\alpha_k(r)\mu_k(m)\).
In the source's commutative algebra
\[
 \mathcal D_e=F\oplus F^e\epsilon\oplus F^e\epsilon'\oplus F\eta,
 \qquad \epsilon_k\epsilon'_j=\delta_{kj}\eta,
\]
all other augmentation products are zero.  There is a ring map
\[
 r+m\longmapsto
 \bar r+\sum_k\mu_k(m)\epsilon_k+
              \sum_k\alpha_k(r)\epsilon'_k+\Phi(m)\eta.
\]
The three product checks are the derivation rule, the displayed
formula for \(\Phi\), and \(\alpha_k(\mathfrak r)=0\).
It supplies an actual representation with every original equation.

Conjugate the odd first-order coefficients to
\(\sum_jc_{kj}\kappa_j\), where \(c_{kj}=\mu_k(m_j)\).
The original skew primitives give an exactly orthogonal conjugator
\(1+\sum_k\epsilon_kY_k\).  Replace the first-order part by the
prescribed \(\rho^1=(1+\sum_j\epsilon_j\kappa_j)\bar\rho\),
send the old \(\epsilon'_k\) coefficients to their sums weighted
by \(c_{kj}\), and keep \(\eta\).  The sole mixed-product check is
\[
 \sum_k\Bigl(\sum_jc_{kj}\kappa_j\Bigr)A_k
       =\sum_j\kappa_j\Bigl(\sum_kc_{kj}A_k\Bigr),
\]
with its reverse-order version.  The original block-compatible
flag makes the odd diagonal zero, so \(\rho^1\) has the original
ordinary characters.  Thus group law, polarization and all
ordinary equations survive.  Explicitly, the relation between old
and new coefficients keeps constants and \(\eta\), sends
\(\epsilon_{\rm old}=c\epsilon_{\rm new}\) and
\(\epsilon'_{\rm new}=c^T\epsilon'_{\rm old}\), and is closed
under products.  Apply it to every ordinary factor, with the same
character transformation.  Its old zero product forces the new
constant, \(\epsilon'\) and mixed \(\eta\) parts to vanish;
the remaining \(\epsilon\) part vanishes by the original
ordinary equation for \(\rho^1\).  No inverse of \(c\) is needed.
Divide by \(\rho^1\), subtract one, and apply the square-zero
local correction above.  This gives the prescribed second-order
lift of \((a_{\alpha_{m_j}})_j\).

\emph{First-order compatibility keeps the local flag.}
Let \(U\subseteq T^g\) be the space of lifting tuples and
\(W=\{\alpha:e_j\alpha\in U\ \forall j\}\).
Condition \(\mathrm K_2\) gives \(\operatorname{codim}U\leq g\).
For an even cocycle \(z\) with lifts
\(\Psi_j=\epsilon'_jz+\eta H_j\), put
\(y=z+\sum_j\epsilon_jH_j\).  Then
\(\epsilon'_jy=\Psi_j\), proving all global conditions.

Locally choose one even primitive \(Y\) and second-order primitives
\(Y_j\).  Their discrepancies
\(Q_{jk}=(Y_j)_{\epsilon'_k}-\delta_{jk}Y\) have upper-triangular
coboundary.  Distinct ordinary characters force \(Q_{jk}\) itself
upper triangular, by descending through its lower diagonals.
With the original flags \(w,w^1\), subtract
\[
 C_j=w^1\left(\sum_k\epsilon'_k w^{-1}Q_{jk}w\right)(w^1)^{-1}
\]
from \(Y_j\).  This remains in the allowed second-order space.
Its coboundary is upper triangular with zero diagonal, so the
local condition is unchanged.  The corrected primitives have
coefficients \(\delta_{jk}Y\); their \(\eta\) coefficients assemble
a first-order primitive \(Y'\) with
\(\epsilon'_jY'=Y_j-C_j\).  This proves the local condition for
\(y\), with no averaging.

A basis of \(W\) gives first-order classes with independent
constant parts.  Adjoin the \(g^2\) classes \(\epsilon_j\kappa_i\).
A boundary relation first vanishes by the tangent injection; its
constant primitive then vanishes by the invariant-skew assertion.
The remaining coefficients vanish by the independence of
\(\kappa_i\).  Condition \(\mathrm K_1\) gives
\[
                         \dim W+g^2\leq D_0.
\]
The case \(g=0\) uses \(y=z\).

\emph{The actual presentation and prime conclusion.}
Apply the established padded construction.  Lift its spanning
forms to \(X_a\in\mathfrak m\) and its \(g^2\) independent
coefficient rows \(c_{ij}\) to \(O'\).  Cotangent duality gives
\[
 \mathfrak m=(X_a)+H,\qquad
             \sum_{j,a}c_{ij,a}X_am_j\in HN.
\]
For the second assertion, apply any \(\Phi\) above, use its
actual lifting tuple and the defining row relation on \(U\),
then use separation by the dual of the finite \(F\)-space \(N/HN\).
This proves the entire integral presentation on the original objects.

For two nonzero blocks, stability implies
\(\mathfrak r\nsubseteq Z\).  At \(Q\), stability gives
\(\mathfrak r\nsubseteq Q\); nonsemisimplicity instead gives
\(\operatorname{Ann}_RN\subseteq Q\) by the exact polarized model.
The local incidence theorem gives \(Z\subseteq Q\).
For one block, \(N=0\), and tangent generation gives
\(O'[[X_1,\ldots,X_{D_0}]]\twoheadrightarrow R\).
The dimension hypothesis forces the inverse image of \(Z\) to be
the zero prime of that regular local domain.  Hence \(Z=0\).
All source quantifiers are retained \cite{SelmerLocalPropagationSource}.
\end{proof}

This closes local propagation at every supplied curve.  It does
not produce a stable prime or a nonzero automorphic support point.
The remaining support assembly must still supply its arithmetic seed
and the original curve chain.  No numerical integral annihilator
bound follows from \(\operatorname{Ann}N\subseteq Q\).
